\section{Resumo}

% Objetivos, Metodologia, Indicação sobre os valores medidos, incluindo a sua precisão e exatidão. Identificar os objetivos atingidos.

\subsection{Fundamentação teórica}
Para a análise de circuitos eletrónicos, são usadas as duas leis de Kirchhoff: a lei dos nós e a lei das malhas:

\paragraph{Lei dos nós}
Um nó não acumula carga, ou seja: atribuindo sinais contrários às correntes que entram e saem de um nó, a soma destas é nula:
\begin{equation*}
	\sum_{i=1}^{n} I_i = 0
\end{equation*}

\paragraph{Lei das malhas}
Atribuindo sinais contrários a tensões que aumentam e tensões que diminuem, a soma destas tensões, ao longo de um caminho fechado é nula:
\begin{equation*}
	\sum_{i=1}^{n} V_i = 0
\end{equation*}

\paragraph{Principio da sobreposicao}
Tendo um circuto elétrico alimentado por fuas fontes de tensão, o efeito combinado em qualquer dos compontentes é a soma dos efeitos provocados por cada uma das fontes de tensão, separadamente.

\subsection{Objetivos}
\begin{itemize}
    \item Usar as várias funções de um multímetro, nomeadamente voltímetro e amperímetro; calibração sensibilidades
	 \item Verificar as leis de Kirchhoff
\end{itemize}
